\section{The Equivalence Theorem}
\label{sec:equivalence}

We now prove that representational and reasoning superposition are instances of the same optimization problem.

\subsection{Unified Loss Function}

\begin{definition}[Co-occurrence Matrix]
\label{def:cooccur}
A \emph{co-occurrence matrix} $P \in [0,1]^{n \times n}$ specifies the probability that objects $i$ and $j$ are simultaneously relevant. We require $P_{ii} = p_i$ (marginal probability) and $P_{ij} = P_{ji}$ (symmetry).
\end{definition}

\begin{definition}[Superposition Loss]
\label{def:super_loss}
For encoding matrix $\mathbf{W} \in \R^{d \times n}$ (where $d$ is the bottleneck dimension) and co-occurrence matrix $P$, define:
\begin{align}
\mathcal{L}_{\text{cap}}(\mathbf{W}) &= \sum_{i=1}^{n} p_i \left(1 - \|\mathbf{w}_i\|^2\right)^2 \label{eq:cap_loss}\\
\mathcal{L}_{\text{int}}(\mathbf{W}, P) &= \sum_{i \neq j} P_{ij} \langle \mathbf{w}_i, \mathbf{w}_j \rangle^2 \label{eq:int_loss}\\
\mathcal{L}(\mathbf{W}, P) &= \mathcal{L}_{\text{cap}}(\mathbf{W}) + \mathcal{L}_{\text{int}}(\mathbf{W}, P) \label{eq:unified_loss}
\end{align}
We call $\mathcal{L}_{\text{cap}}$ the \emph{capacity loss} (penalizing deviation from unit norm) and $\mathcal{L}_{\text{int}}$ the \emph{interference loss} (penalizing non-orthogonality weighted by co-occurrence).
\end{definition}

\begin{theorem}[Representational-Reasoning Equivalence]
\label{thm:equivalence}
Both representational and reasoning superposition minimize the unified loss $\mathcal{L}(\mathbf{W}, P)$ with different co-occurrence matrices:
\begin{enumerate}
    \item[(a)] \textbf{Representational:} Setting $P_{ij} = p_i p_j$ (independent activation), minimizing $\mathcal{L}$ is equivalent to minimizing the expected reconstruction error of \citet{elhage2022toy}.
    \item[(b)] \textbf{Reasoning:} Setting $P_{ij} = \Pr[\text{traces } i, j \text{ both valid}]$, minimizing $\mathcal{L}$ is equivalent to the implicit objective of attention in continuous thought.
\end{enumerate}
\end{theorem}

\subsection{Proof of Part (a): Representational Superposition}

\begin{lemma}
\label{lem:rep_expansion}
For the autoencoder of Section~\ref{sec:rep_superposition} with ReLU activation, under the assumption that interference is small (valid when features are sparse), the expected reconstruction error satisfies:
\begin{equation}
\E\left[\|\mathbf{x} - \hat{\mathbf{x}}\|^2\right] = \sum_{i=1}^n p_i (1 - \|\mathbf{w}_i\|^2)^2 + \sum_{i \neq j} p_i p_j \langle \mathbf{w}_i, \mathbf{w}_j \rangle^2 + O(\epsilon^3)
\end{equation}
where $\epsilon = \max_{i \neq j} |p_i p_j \langle \mathbf{w}_i, \mathbf{w}_j \rangle|$ is the maximum expected interference.
\end{lemma}

\begin{proof}
For input $\mathbf{x}$ with $x_i \in \{0,1\}$ independent Bernoulli$(p_i)$, the pre-activation reconstruction of component $i$ is:
\begin{equation}
z_i = (\mathbf{W}^\top \mathbf{W} \mathbf{x})_i = \sum_{j=1}^n x_j \langle \mathbf{w}_i, \mathbf{w}_j \rangle = x_i \|\mathbf{w}_i\|^2 + \sum_{j \neq i} x_j \langle \mathbf{w}_i, \mathbf{w}_j \rangle
\end{equation}

After ReLU, $\hat{x}_i = \max(0, z_i)$. When $x_i = 1$:
\begin{equation}
z_i = \|\mathbf{w}_i\|^2 + \underbrace{\sum_{j \neq i} x_j \langle \mathbf{w}_i, \mathbf{w}_j \rangle}_{\text{interference term}}
\end{equation}

For sparse features, the interference term is typically small. Expanding to second order:
\begin{align}
\E[(x_i - \hat{x}_i)^2 | x_i = 1] &\approx \E\left[\left(1 - \|\mathbf{w}_i\|^2 - \sum_{j \neq i} x_j \langle \mathbf{w}_i, \mathbf{w}_j \rangle\right)^2\right] \\
&= (1 - \|\mathbf{w}_i\|^2)^2 + \sum_{j \neq i} p_j \langle \mathbf{w}_i, \mathbf{w}_j \rangle^2 + O(\epsilon^3)
\end{align}

When $x_i = 0$, false positives contribute $O(\epsilon^2)$ terms. Summing over all components:
\begin{align}
\E[\|\mathbf{x} - \hat{\mathbf{x}}\|^2] &= \sum_{i=1}^n p_i \E[(x_i - \hat{x}_i)^2 | x_i = 1] + O(\epsilon^2) \\
&= \sum_{i=1}^n p_i (1 - \|\mathbf{w}_i\|^2)^2 + \sum_{i=1}^n \sum_{j \neq i} p_i p_j \langle \mathbf{w}_i, \mathbf{w}_j \rangle^2 + O(\epsilon^3) \\
&= \mathcal{L}_{\text{cap}}(\mathbf{W}) + \mathcal{L}_{\text{int}}(\mathbf{W}, P) + O(\epsilon^3)
\end{align}
where $P_{ij} = p_i p_j$ for $i \neq j$ and $P_{ii} = p_i$.
\end{proof}

\subsection{Proof of Part (b): Reasoning Superposition}

\begin{lemma}
\label{lem:reason_expansion}
For continuous thought with trace embeddings $\mathbf{v}_k$ and attention weights $\alpha_k$, minimizing the next-token prediction loss (cross-entropy) implies minimizing the interference loss $\mathcal{L}_{\text{int}}$.
\end{lemma}

\begin{proof}
From Equation~\ref{eq:thought_decomp}, the thought vector is $\mathbf{h}_t = \sum_k \alpha_k^{(t)} \mathbf{v}_k$. The model predicts the next token by computing logits $\mathbf{z} = \mathbf{W}_U \mathbf{h}_T$ where $\mathbf{W}_U$ is the unembedding matrix.

\textbf{Step 1: Decomposing the logit.}
For the correct answer token $y^*$, let $\mathbf{u}_{y^*}$ be its unembedding vector. Suppose trace $k^*$ leads to the correct answer, i.e., $\mathbf{v}_{k^*}$ is aligned with $\mathbf{u}_{y^*}$. The logit for the correct answer is:
\begin{equation}
z_{y^*} = \langle \mathbf{u}_{y^*}, \mathbf{h}_T \rangle = \alpha_{k^*} \langle \mathbf{u}_{y^*}, \mathbf{v}_{k^*} \rangle + \underbrace{\sum_{j \neq k^*} \alpha_j \langle \mathbf{u}_{y^*}, \mathbf{v}_j \rangle}_{\text{interference noise}}
\end{equation}

\textbf{Step 2: Cross-entropy decomposition.}
The cross-entropy loss is $\mathcal{L}_{\text{CE}} = -\log \text{softmax}(z_{y^*})$. For small interference, Taylor expansion gives:
\begin{align}
\mathcal{L}_{\text{CE}} &\approx \mathcal{L}_{\text{ideal}} + \lambda \cdot \text{Var}[\text{interference}] + O(\epsilon^3) \\
&= \mathcal{L}_{\text{ideal}} + \lambda \sum_{j \neq k^*} \alpha_j^2 \langle \mathbf{u}_{y^*}, \mathbf{v}_j \rangle^2 + O(\epsilon^3)
\end{align}
where $\mathcal{L}_{\text{ideal}}$ is the loss with no interference and $\lambda > 0$ depends on the softmax temperature.

\textbf{Step 3: Connecting to $\mathcal{L}_{\text{int}}$.}
When traces $j$ and $k^*$ are both valid (both could lead to the answer), their embeddings must be distinguishable to avoid the model confusing them. Averaging over problems where different traces are correct:
\begin{equation}
\E[\mathcal{L}_{\text{CE}}] \approx \E[\mathcal{L}_{\text{ideal}}] + \lambda' \sum_{\substack{i, j \in V \\ i \neq j}} \langle \mathbf{v}_i, \mathbf{v}_j \rangle^2 + O(\epsilon^3)
\end{equation}

The second term is precisely $\mathcal{L}_{\text{int}}(\mathbf{W}, P)$ with $P_{ij} = \ind[\text{both valid}]$.

\textbf{Step 4: Capacity loss from normalization.}
For stable training, embeddings must have consistent norms. If $\|\mathbf{v}_k\| \ll 1$, the trace contributes negligibly to $\mathbf{h}_t$; if $\|\mathbf{v}_k\| \gg 1$, gradients become unstable. Layer normalization and weight decay implicitly enforce $\|\mathbf{v}_k\| \approx 1$, contributing the $\mathcal{L}_{\text{cap}}$ term.

Thus, minimizing $\mathcal{L}_{\text{CE}}$ implies minimizing $\mathcal{L}(\mathbf{W}, P) = \mathcal{L}_{\text{cap}} + \mathcal{L}_{\text{int}}$.
\end{proof}

\subsection{Completing the Proof of Theorem~\ref{thm:equivalence}}

\begin{proof}[Proof of Theorem~\ref{thm:equivalence}]
Part (a) follows from Lemma~\ref{lem:rep_expansion}: setting $P_{ij} = p_i p_j$ recovers the expected reconstruction error up to higher-order terms that vanish as sparsity increases.

Part (b) follows from Lemma~\ref{lem:reason_expansion}: setting $P_{ij} = \ind[\text{both valid}]$ recovers the implicit objective of the attention mechanism.

Both are instances of minimizing $\mathcal{L}(\mathbf{W}, P) = \mathcal{L}_{\text{cap}}(\mathbf{W}) + \mathcal{L}_{\text{int}}(\mathbf{W}, P)$, establishing the equivalence.
\end{proof}

\begin{remark}[Role of the Co-occurrence Matrix]
The co-occurrence matrix $P$ captures the \emph{structure} of the superposition problem:
\begin{itemize}
    \item \textbf{Representational:} $P_{ij} = p_i p_j$ reflects independent feature activations. High sparsity (small $p_i$) yields small $P_{ij}$, reducing interference cost.
    \item \textbf{Reasoning:} $P_{ij} = \ind[\text{both valid}]$ reflects the branching structure of valid traces. Problems with few valid paths have sparse $P$, enabling more superposition.
\end{itemize}
The unified framework reveals that both phenomena exploit the same mathematical structure: sparse co-occurrence enables packing more objects into limited dimensions.
\end{remark}

\begin{remark}[Interference as Noise vs.\ Exploration]
\label{rem:interference_exploration}
While the mathematics of superposition is equivalent across settings, the \emph{interpretation} of interference differs:

\paragraph{Representational setting: Interference is pure noise.} In feature encoding, interference $\langle \mathbf{w}_i, \mathbf{w}_j \rangle \neq 0$ is always undesirable---it causes feature $j$'s activation to corrupt the reconstruction of feature $i$. The goal is to minimize interference subject to capacity constraints.

\paragraph{Reasoning setting: Interference enables exploration.} In reasoning superposition, maintaining multiple traces simultaneously is a \emph{feature}, not a bug. The superposition $\mathbf{h}_t = \sum_k \alpha_k \mathbf{v}_k$ allows parallel exploration of the solution space, analogous to beam search or BFS. \citet{zhu2025reasoning} show this enables solving problems that sequential (DFS-like) reasoning cannot.

\paragraph{Resolution: Eventual collapse.} The key insight is that reasoning superposition must \emph{eventually} collapse to a single answer. The loss function $\mathcal{L}_{\text{int}}$ does not penalize superposition per se---it penalizes \emph{indistinguishability} between co-valid traces. When traces $i$ and $j$ are both valid, we need $\langle \mathbf{v}_i, \mathbf{v}_j \rangle \approx 0$ so that the model can later distinguish them and select the correct one. Thus:
\begin{itemize}
    \item \textbf{During exploration:} Multiple traces coexist in superposition, with attention weights $\alpha_k$ distributed across valid options.
    \item \textbf{At decision time:} The gating mechanism (Section~\ref{sec:timespace}) collapses the superposition, concentrating weight on the correct trace: $\alpha_{k^*} \to 1$.
\end{itemize}
Low interference $\langle \mathbf{v}_i, \mathbf{v}_j \rangle \approx 0$ ensures this collapse can be performed accurately---the traces remain distinguishable even while superposed. This parallels quantum computing, where orthogonal states can be maintained in superposition and later measured to yield a definite outcome.
\end{remark}
